\documentclass{article}
\usepackage{amsmath,amssymb,hyperref}
\title{Keeping corridor coverage while the formation moves}
\date{}
\begin{document}
\maketitle
The ordered, speed-limited model from
\href{https://github.com/eternal-november/eos-research-programmes/pull/408}{PR~408}
is held fixed, except that the sensing and radio requirements now
hold throughout the motion from placement $A$ to placement $B$.
Each body stays in its obstacle-safe interval with the prescribed
adjacent clearances $g_i>0$. All sensors have radius $b$;
the same $q$-step radio bound used there enforces connectivity
after $q-1$ subsequent capability losses. The sensing requirement
is $q$-fold pointwise coverage on $[L,U]$. This is continuous
kinematic planning with known endpoints and per-coordinate speed
bounds, not distributed synthesis or motion after a new stopping
event. Both endpoints are assumed feasible.

\paragraph{A protected point.}
Let $C_i$ be the convex configurations with robots $i$ through
$i+q-1$ all inside the sensing band for that point, including the
physical and radio constraints. At any configuration the indices
of covering cells form an interval. Choose any start and finish
witness indices. If they differ, the first times a path reaches
the closed union beyond each successive index cut provide
mandatory adjacent intersections in order. At a first crossing,
closedness gives a witness from each side of the cut: the earlier
side by a convergent sequence of preceding times, or the previous
handoff directly for a coincident crossing. The interval property
then supplies the two adjacent indices. Between these points an
ordinary chord stays in one convex cell and respects the original
per-robot speed bound over the elapsed time. Hence choosing the
transition placements and durations in this fixed cell sequence is
an exact linear program, with no search over routes. Adjacent
transitions may be simultaneous when the same placement suffices.

\paragraph{A protected interval.}
Its feasible set is closed under coordinatewise min and max.
For a point $y$, failure of $q$-fold coverage means that at some
cut $i$, $x_i<y-b$ and $x_{i+q}>y+b$, interpreting indices before
and after the list as negative and positive infinity. If the min
of two feasible placements failed in this way, choose the source
of the lower coordinate at $i$; both sources have their coordinate
$i+q$ too high. That placement already fails. For the max, choose
the source of the higher coordinate at $i+q$; both low coordinates
are too low. This holds for all $y$. Difference constraints and
fixed coordinate intervals are also closed under the two operations.

Let $M=\min(A,B)$. For any path between $A,B$, taking its pointwise
min with $A$ or $B$ gives paths $A$ to $M$ and $M$ to $B$ at no
greater duration, preserving the coordinate speed bounds.
Conversely, pad those two paths to equal time and take their
pointwise max to get a feasible $A$ to $B$ path. For comparable
endpoints $M\le A$, a path clipped between the endpoints can be
made nondecreasing by running coordinate maxima. Every finite
time prefix attains each maximum, which is the join of finitely
many feasible placements. Running maxima respect each speed
bound. Apply this also to $M\le B$ and reverse the first leg.
The exact unrestricted optimum is the maximum of the two
ascending optima, allowing infinity.

\paragraph{A narrow interval has a prescribed event order.}
Suppose $0<U-L<\min_i g_i$. For an ascending leg, use thresholds
$\theta=L-b$ and $\phi=U-b$. The leftmost possible witness for
each end is its first robot not below the corresponding threshold.
These witnesses suffice when the zone is covered: from the
leftmost witness at $L$ to the leftmost at $U$ the ordered common
coverage intervals must overlap all the way. Otherwise any gap
would also be a gap for all later witnesses before $U$.
Each robot crosses the two thresholds at most once for this purpose
(use first arrival at a threshold), and a higher-index robot
passes $\phi$ before a lower-index robot reaches $\theta$.
This follows directly from the clearance exceeding the width.
Retain only the crossings between the known endpoints, with a
crossing at the ending cut allowed.

This determines the entire sequence of witness pairs. On each
stage impose the two end tests and all adjacent overlaps from
PR~408, and impose both threshold groups with weak inequalities
at transitions. Pin the relevant coordinate to the threshold at
each switch, and impose nonnegative coordinate increments on
each segment. This is one linear program per ascending leg,
with at most $2n$ transitions. Necessity follows by taking those
points from a path and using closedness at each cut. Sufficiency
uses a chord in each convex stage; stationary intervals at a
threshold and coincident endpoint changes are harmless. The
ordering is not asserted for a width at least the minimum clearance.
The min/max reduction still holds for wide intervals, but no
polynomial solver for their ascending legs is supplied here.

\paragraph{A tight continuously covered motion.}
For three robots, $q=2$, $b=4$, $[L,U]=[-1,1]$, clearance 4,
radio range 13 and unit speed, the endpoints are
$A=(-12,-3,1)$, $B=(-1,3,12)$. A radio- and separation-feasible
straight move of duration 11 loses coverage. To change which pair
covers the zone, their common windows must touch or overlap:
$x_3-x_1\le8$. Clearance forces equality, so the switch is at
$(s-4,s,s+4)$. Robot 1 needs at least $s+8$ time to get there and
robot 3 needs at least $8-s$ to finish. The lower bound of 16 is
attained via $(-5,-1,3)$ and $(-3,1,5)$ in durations 7, 2, 7.
On the middle chord the two windows meet at the middle robot.
With clearance 5 and range 14, endpoints
$(-12,-3,2)$ and $(-2,3,12)$ remain feasible, but a switch
cannot occur, so there is no continuously covered path.

\paragraph{Matched comparator.}
\href{https://arxiv.org/abs/2205.04422}{Graphs of convex sets}
is the general planning comparison. The direct oracle enumerates
simple cell routes and solves the same speed-constrained chord
problem on each. Revisiting a cell cannot improve an optimum:
replace the intervening path by one chord in that cell.
The restricted construction supplies the route without discrete
search. Numerical comparisons against that exact formulation
test the implementation; they do not replace the geometric proof
or constitute a general planner runtime comparison.
Unlike sensor deployment that only checks the final placement,
this requirement forbids shedding a live sensing duty while
moving. No R6 closure or independently selected I2 claim follows.
\end{document}
